\chapter{Análisis del transporte público}

	La eficiencia y sostenibilidad del transporte público son aspectos críticos en el desarrollo y calidad de vida urbana. Existen múltiples disciplinas tales como la urbanística, la geografía o la sociología que se interesan por abordar este tema desde sus particulares perspectivas. De ahí surge la necesidad de metodologías analíticas que ofrezcan una base rigurosa para diagnosticar, evaluar y planear intervenciones eficaces en los sistemas de transporte.

	Para comenzar, este capítulo introduce algunos de los principales modelos matemáticos y herramientas analíticas que se han empleado en el estudio del transporte público. Estas metodologías permiten medir la eficiencia, la accesibilidad, la vulnerabilidad y otros aspectos estructurales de las redes de transporte.

	De manera particular, reduciremos el enfoque a la aplicación de la teoría de grafos, por su capacidad para representar formalmente las relaciones entre nodos y trayectos. Se explorará cómo esta herramienta se articula con otros elementos como los datos GTFS, y se presentarán casos de estudio que demuestran su utilidad práctica.

    De acuerdo a \cite{rodrigues2020}, el análisis de transporte público se puede segmentar en cuatro niveles, cada uno dependiente del anterior:
    \begin{enumerate}
        \item \textbf{Distancia}: El más esencial, de los cuatro. Consiste en conocer la distancia recorrida entre distintas terminales; el concepto de \textit{distancia} se refiere a cualquier función que cumpla con las propiedades de una norma.
        \item \textbf{Accesibilidad}: Es la medida de qué tan sencillo es alcanzar una cierta locación partiendo de cualquier otra. Es claro que la accesibilidad aquí es una función del lugar y la distancia.
        \item \textbf{Interacciones espaciales}: Es el conjunto de desplazamientos realizados por los usuarios del transporte público. Se toma en cuenta el origen y el destino de los viajes para contrarrestar la oferta con la demanda existente.
        \item \textbf{Modelos de uso del transporte}: Es un marco que "busca evaluar las numerosas relaciones y efectos de retroalimentación entre el transporte y la estructura espacial" \cite{rodrigues2020}.
    \end{enumerate}

    Dependiendo del nivel de abstracción que tengan los datos disponibles, existen múltiples métodos que pueden aplicarse para su estudio. Estos métodos pueden clasificarse de distintas formas según \cite{rodrigues2020}:
        \begin{itemize}
            \item Dependiendo de si son cuantitativos o cualitativos.
            \item Dependiendo de si toman en cuenta a las infraestructuras disponibles.
            \item Dependiendo de si proveen interpolación o extrapolación.
            \item Dependiendo de si el método provee descripción, explicación u optimización.
            \item De acuerdo al nivel de agregación de los datos.
        \end{itemize}

    En el presente trabajo nos enfocaremos en metodologías cuantitativas, concretamente aquellas que ofrecen una descripción del transporte público. En el siguiente capítulo nos extenderemos respecto al tipo de agregación del que dispondremos para nustro modelo propuesto.

    Dentro de los métodos cuantitativos, existen aquellos que parten de una visión multidisciplinaria, como la cartografía, y aquellos que están específicamente enfocados en el transporte y su estudio. Dentro de estos últimos, podemos mencionar a los métodos numéricos enfocados en realizar simulaciones especializadas en el uso del transporte y el suelo. También existe el uso de procesos de optimización que se especializan en responder preguntas tales como dónde ubicar una tienda de retail o una nueva escuela. Finalmente, es común el uso de la teoría de grafos para modelar las redes de transporte público y analizar su estructura. Dentro del presente estudio, nos enfocaremos específicamente en el uso de teoría de grafos para el estudio de redes de transporte público.

    \section{El transporte público como grafo}
        En el capítulo anterior nos propusimos conseguir definir lo que es un grafo, así como sus propiedades principales y varios resultados interesantes en torno a estos. Algunas de las definiciones que adquirimos fue la de recorrido, que nos proporciona una noción de movimiento dentro de un grafo. Asimismo, definimos un tipo de distancia entre dos nodos que nos puede indicar cómo de largo es el recorrido más corto entre estos.

        Nuestra meta en esta sección es mostrar que cualquier red de transporte público puede modelarse como un grafo haciendo uso de toda la teoría vista en el capítulo anterior. Más aún, algunos de los resultados que estudiamos para grafos, también nos dicen cosas muy interesantes del transporte público en tanto tal. Esto hace a los grafos una herramienta valiosísima para modelar esta clase de redes y analizar su comportamiento.

        \subsection{Espacios de representación}
            Un primer reto que tenemos al enfrentar nuestra tarea en esta sección es establecer cuáles son las partes de un sistema de transporte público que nos permitan luego modelarlas como una red. En este sentido, \cite{rodrigues2020} nos dice que, para crear el modelo que queremos, tenemos que partir de las siguientes componentes:

            \begin{itemize}
                \item paradas y estaciones.
                \item intersecciones.
                \item rutas.
            \end{itemize}

            % AGENTE PON AQUÍ UN TIKZ USANDO EL TIKZSET DEL ARCHIVO MAIN.TEX HDP

            Asimismo, nos propone que construyamos un grafo a partir de tres reglas: 1) cada parada o estación, corresponde a un nodo; 2) si es posible ir de una parada o estación a otra a través de una ruta, una arista las une; 3) si dos aristas se intersectan, la intersección corresponde a otro nodo.

            Las reglas 1 y 2 son estándar en la literatura \cite{barraza2021, fielbaum2016, vanferber2009, dingluo2019}, mientras que la regla 3 no está tan presente en otros trabajos, pero es útil en los casos en los que queremos que el grafo resultante del modelo sea también un \textit{grafo planar}.

            \subsubsection{$L\text{-espacio}$}

                Siguiendo estas reglas, hemos creado un $L\text{-espacio}$ o \textit{espacio de infraestructura}, como lo definen \cite{vanferber2009} y después \cite{dingluo2019}. Este modelo es ampliamente utilizado en la literatura \cite{dingluo2019} y suele utilizarse en grafos no ponderados. Es especialmente útil para evaluar la topología de la infraestructura del transporte público. Para una mayor claridad, lo definimos a continuación.

                \begin{definicion}[$L$-espacio]
                    Sea $V$ el conjunto de paradas o estaciones de un sistema de transporte, y sea $E_L$ el conjunto de aristas que representan conexiones físicas directas entre ellas. El grafo del $L$-espacio se define como
                    $$
                        G_L = (V,E_L),
                    $$
                    donde
                    $$
                        (i,j)\in E_L \iff \text{existe un tramo de ruta que conecta directamente las paradas } i \text{ y } j.
                    $$
                \end{definicion}

                En el $L\text{-espacio}$, las aristas representan la existencia de un vínculo físico, no necesariamente de una misma línea o servicio, sino de un segmento compartido de la infraestructura. Esta característica lo convierte en un modelo idóneo para estudiar la \textit{topología estructural} del sistema de transporte, ya que preserva la continuidad espacial de la red, su forma geométrica y sus restricciones físicas \cite{dingluo2019,vanferber2009}. A partir de este modelo pueden calcularse métricas fundamentales como la conectividad global, el diámetro de la red, la longitud media de los caminos, la densidad de nodos o la distribución de grados, las cuales permiten caracterizar la eficiencia, redundancia y resiliencia de la infraestructura ante fallas o bloqueos locales.

                Otra ventaja del $L\text{-espacio}$ es que permite estudiar la relación entre la estructura del transporte y la morfología urbana. En particular, al representar explícitamente las intersecciones, el modelo puede extenderse para incluir como nodos tanto las paradas como los cruces de vías relevantes, de manera que:
                $$
                V = V_p \cup V_i,
                $$
                donde $V_p$ representa el conjunto de paradas y $V_i$ el de intersecciones. En este caso, dos tipos de aristas pueden definirse: las que conectan paradas consecutivas dentro de una misma ruta y las que vinculan intersecciones adyacentes, enriqueciendo el análisis de continuidad espacial del sistema y su integración con la red vial.

                El $L\text{-espacio}$ se corresponde con lo que la literatura denomina un \textit{grafo planar}, es decir, un grafo que puede representarse en el plano sin que sus aristas se crucen. Esta propiedad es especialmente relevante para redes de transporte urbano, donde la estructura topológica se encuentra restringida por la disposición física del territorio y la infraestructura \cite{rodrigues2020}. En contextos metropolitanos complejos, como el de Guadalajara, dicha representación permite relacionar directamente la geometría de las rutas con la accesibilidad territorial y con los patrones de concentración de la demanda \cite{gobjalisco2024,inegi2024,iieg2024}.

                Formalmente, el $L\text{-espacio}$ puede considerarse una proyección de la red bipartita general $G_B=(V\cup R,E_B)$ sobre el conjunto $V$ de paradas, donde dos vértices se conectan si existe una ruta $r\in R$ que las enlaza consecutivamente. Esta proyección conserva la estructura local de conectividad y sirve como base para definir métricas estructurales como:
                $$
                C_i = \frac{2e_i}{k_i(k_i - 1)}, \qquad L = \frac{1}{|V|(|V|-1)} \sum_{i\neq j} d(i,j),
                $$
                donde $C_i$ es el coeficiente de agrupamiento del nodo $i$, $k_i$ su grado y $L$ la longitud media de los caminos en el grafo. Dichos indicadores describen la cohesión topológica de la red, su compacidad y la eficiencia promedio de los desplazamientos posibles a lo largo de la infraestructura \cite{balle1998,diestel2017}.

            % AGENTE PON AQUÍ UN TIKZ USANDO EL TIKZSET DEL ARCHIVO MAIN.TEX HDP

            \subsubsection{$P\text{-espacio}$}

                El $P$-espacio —o \textit{espacio de rutas}— constituye una representación alternativa de la red de transporte público en la cual las relaciones entre paradas no se definen por su conexión física inmediata, como en el $L$-espacio, sino por su pertenencia funcional a una misma ruta o línea de servicio. En este modelo, dos paradas o estaciones se consideran adyacentes si existe al menos una ruta que las conecta de manera directa, sin necesidad de transbordo \cite{dingluo2019,barraza2021}. En otras palabras, las aristas del grafo representan la posibilidad de desplazarse entre dos puntos de la red bajo una misma operación o trayecto, y no la existencia de una infraestructura material que los enlace.

                \begin{definicion}[P-espacio]
                    Sea $R$ el conjunto de rutas o servicios de transporte y, para cada $r\in R$, sea $V_r$ el conjunto de paradas servidas por la ruta $r$. El grafo del $P$-espacio se define como $G_P=(V,E_P)$, donde $V=\bigcup_{r\in R} V_r$ y
                    $$
                        (i,j)\in E_P \iff \exists r\in R \text{ tal que } i,j\in V_r.
                    $$
                    Es decir, dos vértices están conectados si comparten al menos una ruta en común.
                \end{definicion}

                De esta manera, el $P$-espacio modela la red como un conjunto de \textit{cliques} (subgrafos completos) correspondientes a las rutas del sistema, en los cuales todos los vértices están mutuamente conectados. La superposición de estas cliques genera un grafo denso, que revela con claridad la estructura de integración del sistema desde la perspectiva de los desplazamientos posibles sin transbordo. En términos de teoría de grafos, el $P$-espacio puede interpretarse como una proyección del grafo bipartito $G_B=(V\cup R,E_B)$ sobre el conjunto de vértices $V$, donde $(v,r)\in E_B$ indica que la parada $v$ pertenece a la ruta $r$.

                Esta representación permite estudiar propiedades globales y locales que reflejan la \textit{eficiencia funcional} del sistema. A nivel global, el análisis de conectividad, diámetro y densidad de $G_P$ proporciona una medida de cuán integrado está el sistema de transporte sin considerar transbordos, mientras que, a nivel local, las métricas de centralidad —como el grado, la intermediación o la cercanía— indican qué paradas concentran mayor conectividad operacional dentro de sus rutas. En particular, los vértices con alto grado en el $P$-espacio corresponden a paradas con un número elevado de conexiones posibles sin cambio de línea, lo que las convierte en nodos de alta accesibilidad dentro de la estructura de servicio \cite{fielbaum2016,ki2023}.

                El $P$-espacio es, por tanto, una representación orientada a la experiencia del usuario y a la operación del sistema más que a su soporte físico. Mientras que el $L$-espacio permite evaluar la continuidad de la infraestructura, el $P$-espacio permite estudiar la continuidad del viaje. Este enfoque es particularmente relevante cuando se pretende analizar la integración de servicios troncales y alimentadores, o la eficiencia de un rediseño de rutas en términos de reducción de transbordos. En el caso de la Zona Metropolitana de Guadalajara, por ejemplo, la estructura de rutas troncales y complementarias documentada por la Secretaría de Transporte \cite{gobjalisco2024} y los patrones de movilidad registrados por el IIEG \cite{iieg2024} pueden representarse adecuadamente en este modelo, ya que permite detectar redundancias entre líneas, superposiciones de cobertura y cuellos de botella en la conectividad entre servicios.

                Desde un punto de vista analítico, el $P$-espacio también puede ponderarse incorporando atributos asociados a la oferta de transporte, como la frecuencia de paso, la longitud de las rutas o el volumen de pasajeros atendidos. Así, cada arista puede definirse con un peso $w_{ij}$ proporcional al número de rutas que comparten ambas paradas o al flujo estimado de usuarios entre ellas:
                $$
                w_{ij} = f\big(|\{r\in R : i,j\in V_r\}|\big),
                $$
                donde $f$ es una función de ponderación que puede calibrarse según los objetivos del análisis (por ejemplo, para modelar tiempos esperados o demanda efectiva). Este tipo de extensión transforma el grafo en una red ponderada que describe la estructura potencial de movilidad dentro del sistema.

            % AGENTE PON AQUÍ UN TIKZ USANDO EL TIKZSET DEL ARCHIVO MAIN.TEX HDP

            \paragraph{}
            Aunque la gran mayoría de la literatura se centra en estas dos formas de modelar una sistema de transporte público \cite{dingluo2019}, existen otras dos formas que vamos a revisar brevemente, pero que no utilizaremos en el resto de este trabajo.

            \subsubsection*{El $C$-espacio}

                El $C\text{-espacio}$ o \textit{espacio de conexión} amplía el enfoque del $P\text{-espacio}$ al considerar no sólo las paradas servidas por una misma línea, sino las interacciones entre diferentes líneas o modos de transporte que comparten un punto de transferencia. En este modelo, los nodos representan rutas o servicios completos y las aristas se establecen cuando dos rutas tienen al menos una parada en común. Así, el grafo $G_C=(R,E_C)$ está compuesto por el conjunto de rutas $R$ y un conjunto de aristas $E_C$ tal que $(r_i,r_j)\in E_C$ si y sólo si existe $v\in V$ con $v\in V_{r_i}\cap V_{r_j}$.

                El $C\text{-espacio}$ traslada la atención desde la movilidad del usuario individual hacia la \textit{estructura de interconectividad del sistema}. Permite identificar los puntos de transferencia crítica, la redundancia entre líneas y la robustez del sistema frente a perturbaciones localizadas. En la práctica, su análisis se asocia con el estudio de la \textit{resiliencia operacional}, al revelar cómo los flujos pueden redistribuirse entre rutas en caso de fallas o cierres parciales. Esta representación también resulta útil para evaluar el grado de integración modal dentro del transporte metropolitano, al modelar explícitamente las relaciones entre líneas de autobús, tren ligero o BRT bajo un mismo marco topológico \cite{rodrigues2020,diestel2017,balle1998}.

                De manera análoga a los grafos de correlación en otras áreas, el $C\text{-espacio}$ permite aplicar métricas de centralidad sobre rutas —en lugar de paradas—, proporcionando una medida directa de la importancia estratégica de cada línea en la conectividad del conjunto. Así, la topología del $C\text{-espacio}$ complementa la visión local del $L\text{-espacio}$ y la visión funcional del $P\text{-espacio}$, constituyendo una capa intermedia de análisis que describe la forma en que los distintos subsistemas se articulan en un entorno metropolitano.

            \subsubsection*{El $B$-espacio}

                Una última forma de modelar los sistemas de transporte público, son los $B\text{-espacios}$, pero antes de definirlos, debemos definir lo que es un grafo bipartito.

                \begin{definicion}[Grafo bipartito]
                    Sea $G=(V,E)$ un grafo simple. Diremos que $G$ es \textit{bipartito} si existe una partición de su conjunto de vértices $V=V_1\cup V_2$, con $V_1\cap V_2=\emptyset$, tal que toda arista de $E$ une un vértice de $V_1$ con uno de $V_2$; es decir,
                    \[
                        E \subseteq \{(u,v)\mid u\in V_1,\, v\in V_2\}.
                    \]
                \end{definicion}

                El $B\text{-espacio}$, o \textit{espacio bipartito}, constituye una representación más general del sistema de transporte al modelarlo como un grafo bipartito $G_B=(V\cup R, E_B)$, donde un conjunto de vértices corresponde a las paradas $V$ y el otro a las rutas $R$. Una arista $(v,r)\in E_B$ existe si la parada $v$ pertenece a la ruta $r$. A diferencia del $L\text{-espacio}$ y el $P\text{-espacio}$, que surgen de proyecciones de esta estructura bipartita sobre uno de sus conjuntos de vértices, el $B\text{-espacio}$ conserva simultáneamente la relación de pertenencia entre ambos dominios, permitiendo una descripción más completa y neutra del sistema \cite{diestel2017,balle1998}.

                Este modelo resulta útil cuando se busca analizar propiedades de cobertura, redundancia o centralidad conjunta entre paradas y rutas, ya que las medidas calculadas sobre el grafo bipartito reflejan simultáneamente la importancia estructural de las estaciones y de los servicios que las conectan. Además, la estructura $B$-bipartita puede proyectarse de manera natural sobre cualquiera de los dos conjuntos, generando como casos particulares los modelos $L$- y $C$-espacio, según se proyecte sobre paradas o sobre rutas. De esta forma, el $B\text{-espacio}$ actúa como el nivel más general de formalización topológica del sistema de transporte, a partir del cual pueden definirse los demás espacios según el tipo de interacción o análisis requerido.

                Desde el punto de vista metodológico, el $B\text{-espacio}$ constituye la base de una modelación modular y extensible: permite incorporar pesos derivados de frecuencia, capacidad, tiempos de transferencia o flujos de pasajeros, y facilita la integración con modelos de matriz origen–destino o de propagación de perturbaciones \cite{zhao2024,rodrigues2020}. En síntesis, el grafo bipartito $G_B$ ofrece una representación formalmente completa de la red de transporte público, en la que las relaciones estructurales entre infraestructura, operación y movilidad pueden ser exploradas mediante un único marco matemático.